\begin{table}[tbp]
\caption{Observation, measurement, and comparison units.}
\label{tab:ch3-units}
\begingroup\fontsize{9.2}{11.8}\selectfont
\renewcommand{\arraystretch}{1.12}
\begin{tabularx}{\textwidth}{@{}>{\raggedright\arraybackslash}p{23mm}>{\raggedright\arraybackslash}X >{\raggedright\arraybackslash}X@{}}
\toprule
\textbf{Research design} & \textbf{Recorded and measured units} & \textbf{Comparison and interpretation}\\
\midrule
Activity & Events are attributed to accounts and repositories; eligible behaviours form contributor and repository measures. & Distributions, periods, maturity groups and selected projects. Accounts can occur in more than one group.\\[5pt]
OpenRank & Typed nodes and relations form an ecosystem or community graph; the method assigns node scores. & Ranked entities within a declared population; seed sets within a stated diffusion task. Graph boundaries determine the comparison.\\[5pt]
Leaderboard deployment & Public collaboration records are aggregated into project-month outcomes; surveys and interviews concern participants. & Projects and time periods around launch. Repeated monthly records and shared participants need not be independent.\\
\bottomrule
\end{tabularx}
\endgroup
\tablenote{Book synthesis of \citet{xia2022,han2026rank,zhao2024}. Event counts, graph nodes, people and communities are different units; a large number of records does not enlarge the number of independently deployed arrangements.}
\end{table}
